\section{讨论}\label{sec:discussion}

\subsection{重叠窗口引起的随机化单位错配}\label{sec:disc-validity}

本文的理论和模拟共同表明，重叠窗口带来的推断问题超出了“相关样本使有效样本量下降”的标量解释范围。窗口重用首先改变派生特征的联合分布：即使原始观测独立同分布，相邻特征仍可能共享大量底层随机变量，并形成随窗口距离变化的协方差。全排列操作保留经验边际分布，却打散原始顺序中的滞后协方差，因此科学零假设下的“无变化”无法自动提供窗口层面的排列不变性。命题~\ref{prop:counterexample}给出了这一层级错配的有限样本反例，式\eqref{eq:permutation-fpc}则说明排列参考的条件变异只由经验边际量决定。该机制与依赖时间序列中置换检验需要附加有效性条件的既有理论一致\cite{RomanoTirlea2022}。

有效样本量仍可作为特定线性统计量的直观概括。例如，$W=64$、$s=16$、$M=192$ 时，窗口均值的方差比约为 $3.974$。然而，高维最大扫描同时受到方向性、协方差谱和极值运算影响，单一方差比无法确定失校准的方向与幅度。本文逐格结果中同时出现保守、接近名义和反保守状态，也支持这种条件化解释。

\subsection{双预算视角与条件化校准表面}\label{sec:disc-surface}

固定窗口数和固定观测长度回答不同的条件问题。前者保持扫描序列长度 $M$ 不变，步长变化同时改变原始观测长度；后者保持 $T$ 不变，步长变化同时改变窗口数。两种制度均无法单独隔离所有设计因素，但它们提供了互补的预算视角。P1 与 P2 在两个预设数据生成过程的等权平均下得到相近效应量，说明严重失校准并非某一种预算记账方式的偶然产物。相应结论仍应限定于本文的 $W$、步长网格、数据生成过程和表示实现。

校准表面的概念还强调，汇总平均不能替代逐格有效性判断。朴素置换下，BPI 与 MOMENTS 的拒绝率范围分别为 $0.0472$--$0.9756$ 和 $0.0432$--$0.9996$。依赖乘子校准下，相应范围分别为 $0.0056$--$0.1180$ 和 $0.0020$--$0.1412$。将条件间的保守性与反保守性合并，可能得到接近 $0.05$ 的总体均值，同时掩盖局部严重失真。因此，本文将逐格拒绝率、Wilson 区间和预设实用等价带作为主要证据。该分类依赖 Monte Carlo 样本量和等价带宽度，属于操作性有限样本判据。

\subsection{依赖乘子校准的改善与边界}\label{sec:disc-dmb}

M1 在多数重叠配置中大幅降低了 M0 的极端反保守性，说明相关乘子能够向参考分布重新引入部分局部依赖信息。不过，M1 仍同时产生实质保守和实质反保守配置。一般 dependent wild bootstrap 或依赖乘子理论主要针对特定均值、光滑函数或经验过程\cite{Shao2010,BucherKojadinovic2016}；本文管线还包含局部点云构造、持续同调、800 维固定网格映射和候选位置最大化，这些非线性环节使一般渐近结论无法直接转化为统一的有限样本保证。

当前带宽规则也有助于解释非重叠端的保守性。即使 $s=W$，式\eqref{eq:dmb-bandwidth}中的 $\lceil M^{1/3}\rceil$ 仍使 $b>1$；以 P3 的固定 $M=192$ 配置为例，$b=6$，相关乘子不会退化为独立乘子。P3 检出的保守偏差与这一结构相容，但本文没有进一步分解带宽、桥接和高维特征谱各自的贡献。后续研究可针对窗口依赖尺度推导步长适配的乘子协方差，或在原始时间轴上实施块重采样并重新构造全部窗口特征。此类扩展需要新的理论和独立确认实验。

\subsection{表示敏感性与校准有效性}\label{sec:disc-representation}

形状变化和均值平移中的功效反转说明，表示选择决定了观测统计量能够保留何种信号。BPI 在圆形到八字曲线的变化中具有很高的尺寸标准化功效，MOMENTS 在均值平移中具有对应优势。尺寸标准化临界值减少了零假设校准松紧对这一比较的直接影响，因此结果更接近表示层面的信号选择性。与此同时，该比较仍受窗口宽度、表示超参数、变化强度和有限样本设计约束。

表示敏感性与校准有效性应作为两个不同评价维度。信息不足的表示难以通过更准确的校准获得高功效；失真的参考分布也会使信息充分的表示产生错误的显著性判断。两者在有限样本下还可能通过维数、依赖谱和乘子带宽发生交互，本文没有假设其统计独立或严格正交。定位结果进一步显示，观测扫描曲线决定无条件位置误差，校准方法通过拒绝事件影响“检出且定位成功”的联合概率；这两个指标回答的统计问题需要保持区分。

\subsection{实践含义}\label{sec:disc-practice}

对采用重叠窗口的拓扑变化检测，窗口应被视为派生统计单位，而不能因原始观测独立同分布就直接实施窗口级全排列。方法报告至少应明确窗口宽度、步长、原始长度、窗口数和校准单位，并在与实际管线匹配的零假设下逐条件评估 I 类错误率。固定观测预算和固定窗口数对应不同的部署约束，二者的结论不宜互相替代。依赖感知重采样可以作为更合理的参考起点，其带宽和有限样本行为仍需针对具体表示与扫描统计量验证。

\subsection{研究局限}\label{sec:limitations}

本文的结论具有以下边界。
\begin{enumerate}
\item BPI 网格、核宽、尺度上界、同调通道和 MOMENTS 构造均保持固定；本文未系统遍历表示超参数，也未研究特征维数与校准行为之间的敏感性。
\item 零假设仅包含 N-AR 与 N-CIRCLE 两类合成过程，主要结论依赖两个过程的等权设计；更广的非高斯、长记忆、异方差和非平稳零模型尚待检验。
\item 备择仅包括一种形状变化和一种均值平移。形状构造匹配单时点边际均值与协方差，但没有匹配完整滞后联合分布或二阶过程。
\item M1 是冻结的依赖感知参考方法，本文没有为完整 BPI 最大扫描建立新的有限样本或渐近有效性定理，当前带宽规则也未针对每个步长进行最优选择。
\item 逐格校准分类为点态判断，依赖预设等价带和 Monte Carlo 精度，不提供整个设计域上的同时保证。
\item 本稿未纳入自然数据或独立外部数据集，外部有效性主要由合成设计的覆盖范围限定。
\item 本文研究单一变点的最大扫描，尚未处理多变点递归、停止规则及相应的族错误率控制。
\end{enumerate}
